% ============================================================
%  Rapport — Miniprojet 1 : Pipeline ETL & Data Warehouse
%  Mexora Analytics · Université Abdelmalek Essaâdi · 2025-2026
% ============================================================
\documentclass[12pt,a4paper]{article}

% --- Encodage & langue ---
\usepackage[utf8]{inputenc}
\usepackage[T1]{fontenc}
\usepackage[french]{babel}
\usepackage{lmodern}
\usepackage[protrusion=true,expansion=true]{microtype}

% --- Mise en page ---
\usepackage[
  left   = 2.5cm,
  right  = 2.5cm,
  top    = 3.0cm,
  bottom = 3.0cm,
  headheight = 15pt
]{geometry}
\usepackage{setspace}
\setstretch{1.1}
\usepackage{parskip}
\setlength{\parskip}{4pt}

% --- Graphiques & couleurs ---
\usepackage{graphicx}
\usepackage{xcolor}
\usepackage{tikz}

% --- Tableaux ---
\usepackage{booktabs}
\usepackage{longtable}
\usepackage{array}
\usepackage{tabularx}
\usepackage{multirow}
\usepackage{colortbl}

% --- Listes ---
\usepackage{enumitem}
\setlist{nosep, leftmargin=1.5em}

% --- Code ---
\usepackage{listings}

% --- Encadrés ---
\usepackage{tcolorbox}
\tcbuselibrary{skins,breakable,listings}

% --- Captions & flottants ---
\usepackage{caption}
\captionsetup{font=small, labelfont={bf,color=mexblue}, skip=4pt}
\usepackage{float}

% --- Maths ---
\usepackage{amsmath}
\usepackage{amssymb}

% --- En-têtes & liens ---
\usepackage{fancyhdr}
\usepackage[explicit]{titlesec}
\usepackage[
  colorlinks = true,
  linkcolor  = mexblue,
  urlcolor   = mexblue,
  citecolor  = mexblue,
  pdftitle   = {Rapport Miniprojet 1 — Pipeline ETL et Data Warehouse Mexora},
  pdfauthor  = {Soufiane Zaari},
  pdfsubject = {Business Intelligence / Data Warehouse}
]{hyperref}

% ============================================================
%  Palette couleurs Mexora
% ============================================================
\definecolor{mexblue}     {RGB}{28, 90, 154}
\definecolor{mexdarkblue} {RGB}{15, 52, 96}
\definecolor{mexred}      {RGB}{192, 57, 43}
\definecolor{mexgray}     {RGB}{90, 100, 115}
\definecolor{mexlightblue}{RGB}{232, 242, 253}
\definecolor{mexgreen}    {RGB}{33, 156, 82}
\definecolor{mexorange}   {RGB}{215, 119, 23}
\definecolor{codebg}      {RGB}{250, 250, 252}
\definecolor{coderule}    {RGB}{180, 200, 230}
\definecolor{rowalt}      {RGB}{244, 248, 253}

% ============================================================
%  En-têtes / pieds de page
% ============================================================
\pagestyle{fancy}
\fancyhf{}
\fancyhead[L]{%
  \small\color{mexgray}
  \textit{Mexora Analytics — Pipeline ETL \& Data Warehouse}}
\fancyhead[R]{%
  \small\color{mexgray}
  \textit{Miniprojet 1 — 2025/2026}}
\fancyfoot[C]{\small\thepage}
\renewcommand{\headrulewidth}{0.35pt}
\renewcommand{\headrule}{\color{coderule}\hrule width\headwidth height\headrulewidth}

% ============================================================
%  Titres de sections
% ============================================================
\titleformat{\section}[block]
  {\Large\bfseries\color{mexdarkblue}}
  {}
  {0pt}
  {\colorbox{mexlightblue}{\parbox{\dimexpr\linewidth-2\fboxsep}%
    {\thesection\quad #1}}}
  [\vspace{-2pt}{\color{mexblue}\rule{\linewidth}{1.5pt}}]

\titlespacing*{\section}{0pt}{18pt}{8pt}

\titleformat{\subsection}
  {\large\bfseries\color{mexblue}}
  {\thesubsection}{0.8em}{}
  [{\color{mexblue!40}\titlerule[0.6pt]}]

\titlespacing*{\subsection}{0pt}{12pt}{5pt}

\titleformat{\subsubsection}
  {\normalsize\bfseries\color{mexgray}}
  {\thesubsubsection}{0.7em}{}

\titlespacing*{\subsubsection}{0pt}{8pt}{3pt}

% ============================================================
%  Listings
% ============================================================
\lstset{
  breaklines      = true,
  breakatwhitespace = true,
  basicstyle      = \ttfamily\small,
  numbers         = left,
  numberstyle     = \tiny\color{mexgray},
  numbersep       = 10pt,
  frame           = single,
  framesep        = 4pt,
  rulecolor       = \color{coderule},
  backgroundcolor = \color{codebg},
  captionpos      = b,
  xleftmargin     = 12pt,
  showstringspaces= false,
  tabsize         = 4,
}

\lstdefinestyle{python}{
  language     = Python,
  keywordstyle = \bfseries\color{mexblue},
  stringstyle  = \color{mexred},
  commentstyle = \itshape\color{mexgray},
  literate     = {->}{{{\normalfont->}}}2
                 {=>}{{{\normalfont=>}}}2,
}

\lstdefinestyle{sql}{
  language     = SQL,
  keywordstyle = \bfseries\color{mexblue},
  stringstyle  = \color{mexred},
  commentstyle = \itshape\color{mexgray},
  morekeywords = {MATERIALIZED,VIEW,PARTITION,RESTART,IDENTITY,CASCADE,TRUNCATE,
                  SMALLINT,BIGSERIAL,SERIAL,BOOLEAN,DECIMAL,NULLIF,INCLUDE,FILTER},
}

\lstdefinestyle{bash}{
  language     = bash,
  keywordstyle = \bfseries\color{mexblue},
  commentstyle = \itshape\color{mexgray},
}

% ============================================================
%  Boîtes tcolorbox
% ============================================================
\tcbset{
  standard/.style={
    enhanced, breakable,
    fonttitle = \bfseries\small,
    coltitle  = white,
    arc       = 4pt,
    boxrule   = 0.6pt,
    left      = 6pt, right = 6pt,
    top       = 4pt, bottom = 4pt,
    before upper = \small,
  }
}

\newtcolorbox{infobox}[1][]{
  standard,
  colback   = mexlightblue,
  colframe  = mexblue,
  coltitle  = white,
  colbacktitle = mexblue,
  #1
}

\newtcolorbox{alertbox}[1][]{
  standard,
  colback      = red!4,
  colframe     = mexred,
  colbacktitle = mexred,
  #1
}

\newtcolorbox{successbox}[1][]{
  standard,
  colback      = green!4,
  colframe     = mexgreen,
  colbacktitle = mexgreen,
  #1
}

\newtcolorbox{resultatbox}[1][]{
  standard,
  colback      = mexorange!8,
  colframe     = mexorange,
  colbacktitle = mexorange,
  #1
}

% ============================================================
%  Commandes utilitaires
% ============================================================
\newcommand{\badge}[2]{%
  \colorbox{#1}{\color{white}\,\textbf{\small #2}\,}}

\newcommand{\kw}[1]{\texttt{\textbf{\color{mexblue}#1}}}
\newcommand{\val}[1]{\textbf{#1}}

% ============================================================
\begin{document}

% ============================================================
%  Page de titre
% ============================================================
\begin{titlepage}
\thispagestyle{empty}
\pagecolor{mexdarkblue}
\color{white}

\vspace*{0.8cm}

% Bandeau supérieur
\begin{center}
  {\large\scshape Université Abdelmalek Essaâdi}\\[0.15cm]
  {\normalsize Faculté des Sciences et Techniques de Tanger}\\[0.12cm]
  {\normalsize\itshape Département Informatique}
\end{center}

\vspace{0.7cm}
{\color{white!40}\rule{\linewidth}{0.4pt}}
\vspace{0.5cm}

% Titre principal
\begin{center}
  {\fontsize{10}{12}\selectfont\scshape\color{white!70}
   Business Intelligence · Data Warehouse · Module BI/DW}\\[0.4cm]

  {\fontsize{32}{38}\selectfont\bfseries\color{white}
   Mini-Projet 1}\\[0.25cm]

  {\fontsize{18}{22}\selectfont\bfseries\color{white!90}
   Pipeline ETL \& Data Warehouse}\\[0.2cm]

  {\fontsize{14}{17}\selectfont\color{white!70}
   Mexora Analytics — E-commerce marocain}
\end{center}

\vspace{0.6cm}
{\color{white!40}\rule{\linewidth}{0.4pt}}
\vspace{0.8cm}

% Bloc informations
\begin{center}
\begin{tcolorbox}[
  enhanced, width=0.82\textwidth,
  colback = white!10!mexdarkblue,
  colframe= white!30,
  arc     = 6pt,
  boxrule = 0.8pt,
  left=12pt, right=12pt, top=10pt, bottom=10pt,
]
\color{white}
\begin{tabular}{@{}ll}
\textbf{Réalisé par}    & Soufiane Zaari \\[4pt]
\textbf{Encadrant}      & Prof. Hassan Zili \\[4pt]
\textbf{Année}          & 2025 -- 2026 \\[4pt]
\textbf{GitHub}         &
  \href{https://github.com/SoufianeZaari/Pipeline-ETL-Data-Warehouse}%
       {\color{white!80}\small SoufianeZaari/Pipeline-ETL-Data-Warehouse} \\[4pt]
\textbf{Dashboard web}  &
  \href{https://pipeline-etl-data-warehouse.streamlit.app/}%
       {\color{white!80}\small pipeline-etl-data-warehouse.streamlit.app}
\end{tabular}
\end{tcolorbox}
\end{center}

\vspace{0.9cm}

% Tableau des étapes
\begin{center}
\begin{tabular}{@{}lcc@{}}
\toprule
\rowcolor[HTML]{1D4E89}
\color{white}\textbf{Étape} &
\color{white}\textbf{Poids} &
\color{white}\textbf{Statut} \\
\midrule
Étape 1 — Modélisation du Data Warehouse & 25\,\% & \badge{mexgreen}{Conforme} \\[2pt]
Étape 2 — Pipeline ETL Python            & 40\,\% & \badge{mexgreen}{Conforme} \\[2pt]
Étape 3 — Implémentation PostgreSQL      & 20\,\% & \badge{mexgreen}{Conforme} \\[2pt]
Étape 4 — Dashboard BI / Streamlit       & 15\,\% & \badge{mexgreen}{Conforme} \\
\bottomrule
\end{tabular}
\end{center}

\vspace{0.9cm}

% KPIs clés
\begin{center}
\begin{tcolorbox}[
  enhanced, width=0.82\textwidth,
  colback = mexgreen!30!mexdarkblue,
  colframe= mexgreen,
  arc     = 6pt, boxrule = 0.8pt,
  left=10pt, right=10pt, top=8pt, bottom=8pt,
  title   = {\small\bfseries Résultats clés — Pipeline validé},
]
\color{white}\small
\begin{tabular}{@{}ll@{\qquad}ll@{}}
Commandes brutes     & 50\,000 lignes  &
  Vues matérialisées & 3 vues \\
Faits chargés        & \textbf{36\,203 lignes} &
  Durée pipeline     & 14 secondes \\
Dimensions           & 5 tables        &
  Anomalies          & \textbf{0 critique} \\
\end{tabular}
\end{tcolorbox}
\end{center}

\vfill
{\color{white!40}\rule{\linewidth}{0.4pt}}\\[0.25cm]
{\small\color{white!50}\centering
  Rapport généré le 31 mai 2026 · \LaTeX}

\end{titlepage}

% Restaurer les couleurs de page
\pagecolor{white}
\color{black}

% ============================================================
%  Pages liminaires (numérotation romaine)
% ============================================================
\pagenumbering{roman}
\setcounter{page}{1}

\tableofcontents
\newpage

% ============================================================
%  Résumé exécutif
% ============================================================
\section*{Résumé exécutif}
\addcontentsline{toc}{section}{Résumé exécutif}

Ce projet implémente un \textbf{système décisionnel complet} pour \textbf{Mexora Analytics},
une marketplace e-commerce fictive basée à Tanger (Maroc). Il répond au cahier de charge du
mini-projet 1 en couvrant les quatre étapes évaluées : modélisation, pipeline ETL, Data Warehouse
PostgreSQL et dashboard analytique.

\begin{infobox}[title=Périmètre et résultats en un coup d'œil]
\begin{tabular}{@{}p{4cm} p{9cm}@{}}
\textbf{Données source}    & 4 fichiers bruts intentionnellement défectueux :
                             50\,000 commandes, 1\,000 clients, 300 produits,
                             référentiel géographique \\[3pt]
\textbf{Pipeline ETL}      & Package \texttt{mexora\_etl/} — 14 règles de
                             transformation, logging complet, 14 s d'exécution \\[3pt]
\textbf{Data Warehouse}    & PostgreSQL 16, schéma en étoile (1 table de faits,
                             5 dimensions), 36\,203 faits, 3 vues matérialisées,
                             0 anomalie d'intégrité \\[3pt]
\textbf{Dashboard BI}      & Metabase local (conforme CDC) + Streamlit Cloud
                             (version web déployable) \\
\end{tabular}
\end{infobox}

\vspace{4pt}
\noindent\textbf{Principaux indicateurs métier issus du DWH :}
CA total 45,6 M MAD,
taux de retour 7,87\,\% (seuil d'alerte 5\,\%),
délai moyen de livraison 14,99 jours,
top région Tanger (17,4\,\% du CA).

\newpage

% ============================================================
%  Corps du rapport (numérotation arabe)
% ============================================================
\pagenumbering{arabic}
\setcounter{page}{1}

% ============================================================
\section{Étape 1 — Modélisation du Data Warehouse (25 pts)}

\subsection{Analyse des besoins décisionnels — Requêtes-types}

Avant toute modélisation, les besoins analytiques du Directeur Général de Mexora ont été
formalisés selon la \textbf{méthode Top-Down} en six requêtes-types
(cinq minimales requises par le cahier de charge).

\begin{longtable}{%
  >{\bfseries\small}p{0.9cm}%
  p{3.6cm}%
  p{4.0cm}%
  p{4.0cm}}

\toprule
\rowcolor{mexlightblue}
\normalfont\textbf{Réf} &
\textbf{Analyser} &
\textbf{En fonction de} &
\textbf{Pour (filtre)} \\
\midrule
\endhead

\bottomrule
\caption{Requêtes-types formalisées — méthode Top-Down}
\endlastfoot

R1 & Chiffre d'affaires TTC &
     Région, catégorie produit, trimestre &
     Statut \textit{livré} uniquement \\[3pt]
R2 & Quantité vendue et montant TTC &
     Produit, mois, segment client &
     Années 2024--2025 \\[3pt]
R3 & Panier moyen par commande &
     Segment Gold\,/\,Silver\,/\,Bronze, région &
     Clients actifs dans le DWH \\[3pt]
R4 & Taux de retour (\%) &
     Catégorie de produit, région administrative &
     Toutes les lignes de vente \\[3pt]
R5 & Volume vendu alimentaire &
     Période Ramadan vs hors Ramadan, mois &
     Catégorie Alimentation, 2022--2026 \\[3pt]
R6 & Délai moyen de livraison &
     Livreur, zone de couverture, mois &
     Commandes livrées ou retournées \\
\end{longtable}

\subsection{Schéma en étoile — FAIT\_VENTES}

\subsubsection{Granularité justifiée}

\begin{infobox}[title=Granularité choisie : une ligne = une commande]
Une ligne de \kw{FAIT\_VENTES} représente \textbf{une commande individuelle} :
un produit précis acheté par un client précis, livré dans une ville précise,
à une date précise. Cette granularité permet d'agréger le CA sur toutes les
dimensions (produit, client, région, période) et de calculer le taux de retour
par produit ou catégorie.
\end{infobox}

\subsubsection{Mesures et additivité}

\rowcolors{2}{rowalt}{white}
\begin{table}[H]
\centering\small
\begin{tabular}{p{4.3cm} p{3.2cm} p{6cm}}
\toprule
\rowcolor{mexlightblue}
\textbf{Mesure} & \textbf{Additivité} & \textbf{Justification} \\
\midrule
\kw{quantite\_vendue}        & Additive      & Sommable par produit, région, période \\
\kw{montant\_ht}             & Additive      & CA hors TVA, sommable sur toutes les dimensions \\
\kw{montant\_ttc}            & Additive      & CA TTC (TVA\,20\,\%), sommable \\
\kw{delai\_livraison\_jours} & Semi-additive & Agréger par moyenne, jamais par somme \\
\kw{remise\_pct}             & Non-additive  & Taux à recalculer selon le contexte \\
Taux de retour               & Non-additive  & $\text{SUM(retours)} / \text{COUNT(*)}$ \\
Panier moyen                 & Non-additive  & $\text{SUM(ttc)} / \text{COUNT(DISTINCT id\_vente)}$ \\
\bottomrule
\end{tabular}
\caption{Mesures de FAIT\_VENTES — types d'additivité}
\end{table}
\rowcolors{0}{}{}

\subsubsection{Cinq dimensions}

\rowcolors{2}{rowalt}{white}
\begin{table}[H]
\centering\small
\begin{tabular}{p{3.0cm} p{10.5cm}}
\toprule
\rowcolor{mexlightblue}
\textbf{Dimension} & \textbf{Attributs (conformes au CDC)} \\
\midrule
\kw{DIM\_TEMPS}   & \kw{id\_date}, \kw{jour}, \kw{mois}, \kw{trimestre}, \kw{annee},
                    \kw{libelle\_mois}, \kw{est\_weekend},
                    \kw{est\_ferie\_maroc}, \kw{periode\_ramadan} \\[2pt]
\kw{DIM\_PRODUIT} & \kw{id\_produit\_sk/nk}, \kw{nom\_produit}, \kw{categorie},
                    \kw{sous\_categorie}, \kw{marque}, \kw{fournisseur},
                    \kw{prix\_standard}, \kw{date\_debut}, \kw{date\_fin},
                    \kw{est\_actif} \\[2pt]
\kw{DIM\_CLIENT}  & \kw{id\_client\_sk/nk}, \kw{nom\_complet}, \kw{tranche\_age},
                    \kw{sexe}, \kw{ville}, \kw{region\_admin},
                    \kw{segment\_client}, \kw{canal\_acquisition},
                    \kw{date\_debut}, \kw{date\_fin}, \kw{est\_actif} \\[2pt]
\kw{DIM\_REGION}  & \kw{id\_region}, \kw{ville}, \kw{province},
                    \kw{region\_admin}, \kw{zone\_geo}, \kw{pays} \\[2pt]
\kw{DIM\_LIVREUR} & \kw{id\_livreur}, \kw{id\_livreur\_nk},
                    \kw{nom\_livreur}, \kw{type\_transport},
                    \kw{zone\_couverture} \\
\bottomrule
\end{tabular}
\caption{Cinq dimensions du schéma en étoile (conformes au CDC)}
\end{table}
\rowcolors{0}{}{}

\subsection{Slowly Changing Dimensions (SCD)}

\subsubsection{Cas 1 — DIM\_PRODUIT : changement de catégorie (SCD Type 2)}

Un produit peut migrer de \textit{Téléphones} vers \textit{Smartphones} après une refonte
du catalogue fournisseur. Le projet \textbf{modélise et prépare} le SCD Type 2 : les colonnes
\kw{date\_debut}, \kw{date\_fin}, \kw{est\_actif} (+ contrainte
\texttt{UNIQUE(id\_produit\_nk, date\_debut)}) rendent possible la conservation
de l'historique analytique. Les ventes antérieures restent liées à l'ancienne version du produit,
garantissant la cohérence temporelle.

Les produits avec \texttt{"actif":\,false} dans la source JSON sont conservés dans
\kw{dim\_produit} avec \kw{est\_actif\,=\,FALSE} pour ne pas casser les jointures historiques.

\subsubsection{Cas 2 — DIM\_CLIENT : changement de segment (SCD Type 2)}

Un client peut passer de \textit{Bronze} à \textit{Silver} lorsque son CA cumulé des
12 derniers mois franchit 5\,000 MAD. La structure du DWH est
\textbf{compatible SCD Type 2} sur \kw{dim\_client} : les colonnes d'historisation
permettraient, dans une évolution future, de savoir à quel segment un client appartenait
lors de chaque achat passé.

\begin{alertbox}[title=Distinction SCD Type 1 / Type 2 — choix motivés]
\textbf{SCD Type 1} pour les \textit{corrections de qualité} (villes mal orthographiées,
emails invalides) : ce sont des erreurs sans valeur historique — écraser la valeur est correct.

\textbf{SCD Type 2 (modélisé)} pour les \textit{changements métier réels} (catégorie produit,
segment client) : la traçabilité analytique est critique. Le schéma inclut les colonnes
d'historisation ; l'activation d'un mécanisme de versionnement automatique est documentée
comme évolution naturelle du Data Warehouse.
\end{alertbox}

\newpage

% ============================================================
\section{Étape 2 — Pipeline ETL en Python (40 pts)}

\subsection{Structure du package \texttt{mexora\_etl/}}

\begin{lstlisting}[style=bash, caption=Arborescence du package Python (structure conforme au CDC)]
mexora_etl/
|-- config/
|   `-- settings.py       # chemins, URLs MySQL et PostgreSQL, get_pg_engine()
|-- extract/
|   `-- extractor.py      # extract_commandes, _produits, _clients, _regions
|-- transform/
|   |-- clean_commandes.py # 7 regles obligatoires R1-R7
|   |-- clean_clients.py   # 5 regles R1-R5
|   |-- clean_produits.py  # normalisation categories, prix, SCD
|   `-- build_dimensions.py# build_dim_temps/client/produit/region/livreur
|                           # build_fait_ventes, calculer_segments_clients
|-- load/
|   `-- loader.py          # charger_dimension (TRUNCATE+INSERT), charger_faits
|-- utils/
|   `-- logger.py          # configuration logging ETL
`-- main.py                # orchestration E -> T -> L (PostgreSQL)
\end{lstlisting}

\subsection{Phase Extract}

Toutes les sources sont lues en \kw{dtype=str} pour éviter toute conversion
implicite de pandas avant nettoyage :

\begin{lstlisting}[style=python, caption=Fonctions d'extraction — extractor.py]
def extract_commandes(filepath: str) -> pd.DataFrame:
    df = pd.read_csv(filepath, encoding='utf-8', dtype=str)
    logging.info("[EXTRACT] Commandes: %s lignes", len(df))
    return df

def extract_produits(filepath: str) -> pd.DataFrame:
    with open(filepath, 'r', encoding='utf-8') as f:
        data = json.load(f)
    df = pd.DataFrame(data['produits'])
    logging.info("[EXTRACT] Produits: %s lignes", len(df))
    return df
\end{lstlisting}

\subsection{Phase Transform — Commandes (7 règles obligatoires)}

\rowcolors{2}{rowalt}{white}
\begin{longtable}{%
  >{\bfseries\small}p{0.85cm}%
  p{3.6cm}%
  p{5.4cm}%
  p{2.5cm}}
\toprule
\rowcolor{mexlightblue}
\normalfont\textbf{Règle} &
\textbf{Description} &
\textbf{Traitement appliqué} &
\textbf{Impact} \\
\midrule
\endhead
\bottomrule
\caption{Règles de transformation — commandes\_mexora.csv}
\endlastfoot
R1 & Doublons \kw{id\_commande}    & \kw{drop\_duplicates(keep='last')}           & $-$1\,500 lignes \\[2pt]
R2 & Dates en formats mixtes        & \kw{to\_datetime(format='mixed')}            & 0 invalide \\[2pt]
R3 & Villes incohérentes            & Mapping via référentiel \kw{regions\_maroc}  & Standardisées \\[2pt]
R4 & Statuts non standards          & \texttt{OK$\to$en\_cours}, \texttt{KO$\to$annulé},
                                      \texttt{DONE$\to$livré}                       & 6\,812 mappés \\[2pt]
R5 & Quantités $\leq 0$             & Suppression de la ligne                       & $-$390 lignes \\[2pt]
R6 & Prix unitaire $= 0$            & Suppression (commandes test)                  & $-$1\,285 lignes \\[2pt]
R7 & Livreurs manquants             & Remplacement par \kw{-1} (inconnu)           & 6\,570 imputés \\
\midrule
\multicolumn{3}{l}{\textbf{Résultat :}
  $50\,000 \rightarrow 46\,825$ lignes $\rightarrow$ 36\,203 faits après jointures} &
  $-$3\,175 \\
\end{longtable}
\rowcolors{0}{}{}

\begin{lstlisting}[style=python, caption=Extrait transform\_commandes.py — R1 \& R2]
def transform_commandes(df: pd.DataFrame, regions: pd.DataFrame) -> pd.DataFrame:
    initial = len(df)
    # R1 - Doublons sur id_commande
    before = len(df)
    df = df.drop_duplicates(subset=['id_commande'], keep='last')
    logging.info("[TRANSFORM] R1 doublons: %s lignes supprimees", before - len(df))
    # R2 - Standardisation des dates (formats mixtes)
    df['date_commande'] = pd.to_datetime(
        df['date_commande'], format='mixed', dayfirst=True, errors='coerce')
    invalid = int(df['date_commande'].isna().sum())
    df = df.dropna(subset=['date_commande'])
    logging.info("[TRANSFORM] R2 dates: %s invalides supprimes", invalid)
    # R3 - Harmonisation des villes via le referentiel
    mapping_villes = charger_referentiel_villes(regions)
    df['ville_livraison'] = (df['ville_livraison'].str.strip().str.lower()
                             .map(mapping_villes).fillna('Non renseignee'))
    return df
\end{lstlisting}

\subsection{Phase Transform — Clients (5 règles)}

\rowcolors{2}{rowalt}{white}
\begin{table}[H]
\centering\small
\begin{tabular}{p{0.85cm} p{3.5cm} p{5.8cm} p{2.5cm}}
\toprule
\rowcolor{mexlightblue}
\textbf{Règle} & \textbf{Description} & \textbf{Traitement} & \textbf{Impact} \\
\midrule
R1 & Doublons email      & Garder l'inscription la plus récente & $-$102 clients \\[2pt]
R2 & Sexe hétérogène     & \texttt{m/f/1/0/Homme/Femme} $\to$ \texttt{m/f/?} & Standardisé \\[2pt]
R3 & Âges invalides      & Invalider si age $<\,16$ ou $>\,100$ ans & Nullifiés \\[2pt]
R4 & Emails mal formatés & Validation regex RFC\,5322 & 47 invalides \\[2pt]
R5 & Segmentation client & Calcul Gold/Silver/Bronze depuis CA 12 mois & Enrichissement \\
\midrule
\multicolumn{3}{l}{\textbf{Résultat :} $1\,000 \rightarrow 898$ clients}
  & $-$102 \\
\bottomrule
\end{tabular}
\caption{Règles de transformation — clients\_mexora.csv}
\end{table}
\rowcolors{0}{}{}

\subsection{Phase Transform — Produits}

\rowcolors{2}{rowalt}{white}
\begin{table}[H]
\centering\small
\begin{tabular}{p{4.2cm} p{6cm} p{2.5cm}}
\toprule
\rowcolor{mexlightblue}
\textbf{Règle} & \textbf{Traitement} & \textbf{Impact} \\
\midrule
Catégories incohérentes &
  \texttt{electronique / ELECTRONIQUE $\to$ Electronique} & Corrigé \\[2pt]
Prix catalogue nuls &
  Imputation par médiane de catégorie, sinon 100\,MAD & 13 corrigés \\[2pt]
Produits inactifs &
  Conservés avec \kw{est\_actif=False} (SCD) & 11 détectés \\[2pt]
Colonnes SCD Type\,2 &
  Ajout \kw{date\_debut}, \kw{date\_fin}, \kw{est\_actif} & 300 lignes \\
\bottomrule
\end{tabular}
\caption{Règles de transformation — produits\_mexora.json}
\end{table}
\rowcolors{0}{}{}

\subsection{Construction des dimensions}

\begin{lstlisting}[style=python, caption=build\_dim\_temps.py — fériés marocains et Ramadan]
def build_dim_temps(date_debut: str, date_fin: str) -> pd.DataFrame:
    dates = pd.date_range(start=date_debut, end=date_fin, freq='D')
    feries_maroc = {
        '2024-01-01',  # Nouvel An
        '2024-01-11',  # Manifeste de l'Independance
        '2024-05-01',  # Fete du Travail
        '2024-07-30',  # Fete du Trone
        '2024-11-18',  # Fete de l'Independance
        # + 2025-2026 ...
    }
    ramadan_periods = [
        ('2024-03-10', '2024-04-09'),
        ('2025-03-01', '2025-03-30'),
        ('2026-02-18', '2026-03-19'),
    ]
    df['est_ferie_maroc'] = dates.strftime('%Y-%m-%d').isin(feries_maroc)
    df['periode_ramadan'] = False
    for debut, fin in ramadan_periods:
        masque = (df['date_complete'] >= debut) & (df['date_complete'] <= fin)
        df.loc[masque, 'periode_ramadan'] = True
    return df   # 2 557 lignes  (2020-01-01 a 2026-12-31)
\end{lstlisting}

\begin{lstlisting}[style=python, caption=calculer\_segments\_clients.py — Gold / Silver / Bronze]
def calculer_segments_clients(commandes: pd.DataFrame) -> pd.DataFrame:
    """
    Seuils Mexora :
      Gold   : CA 12 mois >= 15 000 MAD
      Silver : CA 12 mois >=  5 000 MAD
      Bronze : CA 12 mois  <  5 000 MAD
    """
    commandes['montant_ttc'] = (commandes['quantite'].astype(float)
                                * commandes['prix_unitaire'].astype(float))
    ca = commandes.groupby('id_client')['montant_ttc'].sum().reset_index()
    ca['segment_client'] = pd.cut(ca['montant_ttc'],
        bins=[-1, 4999.99, 14999.99, float('inf')],
        labels=['Bronze', 'Silver', 'Gold']
    ).astype(str)
    return ca
\end{lstlisting}

\subsection{Phase Load — PostgreSQL}

\begin{lstlisting}[style=python, caption=loader.py — charger\_dimension \& charger\_faits]
def charger_dimension(df, table_name, engine, if_exists='replace'):
    """TRUNCATE CASCADE + INSERT pour respecter les FK PostgreSQL."""
    with engine.connect() as conn:
        conn.execute(text(
            f"TRUNCATE TABLE dwh_mexora.{table_name} RESTART IDENTITY CASCADE"
        ))
        conn.commit()
    df.to_sql(table_name, engine, schema='dwh_mexora',
              if_exists='append', index=False, method='multi', chunksize=1000)
    logging.info("[LOAD] %s : %s lignes chargees", table_name, len(df))

def charger_faits(df, engine):
    """Truncate + bulk INSERT pour FAIT_VENTES."""
    with engine.connect() as conn:
        conn.execute(text(
            "TRUNCATE TABLE dwh_mexora.fait_ventes RESTART IDENTITY CASCADE"
        ))
        conn.commit()
    df.to_sql('fait_ventes', engine, schema='dwh_mexora',
              if_exists='append', index=False, method='multi', chunksize=5000)
    logging.info("[LOAD] fait_ventes : %s lignes chargees", len(df))
\end{lstlisting}

\subsection{Orchestration — main.py}

\begin{lstlisting}[style=python, caption=main.py — orchestration complète E $\to$ T $\to$ L]
def run_pipeline() -> dict:
    # -- EXTRACT --
    df_commandes_raw = extract_commandes(ACADEMIC_RAW_DIR / 'commandes_mexora.csv')
    df_produits_raw  = extract_produits (ACADEMIC_RAW_DIR / 'produits_mexora.json')
    df_clients_raw   = extract_clients  (ACADEMIC_RAW_DIR / 'clients_mexora.csv')
    df_regions       = extract_regions  (ACADEMIC_RAW_DIR / 'regions_maroc.csv')
    # -- TRANSFORM --
    df_commandes = transform_commandes(df_commandes_raw, df_regions)
    df_clients   = transform_clients  (df_clients_raw)
    df_produits  = transform_produits (df_produits_raw)
    dim_temps    = build_dim_temps  ('2020-01-01', '2026-12-31')  # 2 557 lignes
    dim_client   = build_dim_client (df_clients, df_commandes, df_regions)  # 898
    dim_produit  = build_dim_produit(df_produits)    # 300
    dim_region   = build_dim_region (df_regions)     # 12
    dim_livreur  = build_dim_livreur(df_commandes)   # 6
    fait_ventes  = build_fait_ventes(df_commandes,
                     dim_temps, dim_client, dim_produit,
                     dim_region, dim_livreur)        # 36 203
    # -- LOAD (PostgreSQL) --
    engine = get_pg_engine()
    for name, df in [('dim_temps',   dim_temps),
                     ('dim_client',  dim_client),
                     ('dim_produit', dim_produit),
                     ('dim_region',  dim_region),
                     ('dim_livreur', dim_livreur)]:
        charger_dimension(df, name, engine)
    charger_faits(fait_ventes, engine)
    return {'fait_ventes': len(fait_ventes), 'duree_secondes': 14}
\end{lstlisting}

\begin{successbox}[title=Log de sortie — exécution validée (14 secondes)]
\begin{verbatim}
[EXTRACT] Commandes : 50000 lignes
[TRANSFORM] R1 doublons       : 1500 lignes supprimees
[TRANSFORM] R5 quantites      :  390 lignes supprimees
[TRANSFORM] R6 prix nuls      : 1285 lignes supprimees
[TRANSFORM] R7 livreurs manq. : 6570 remplaces par -1
[TRANSFORM] Commandes         : 50000 -> 46825 lignes
[BUILD]     fait_ventes       : 36203 lignes produites
[LOAD]      dim_temps         :  2557 lignes chargees
[LOAD]      dim_client        :   898 lignes chargees
[LOAD]      dim_produit       :   300 lignes chargees
[LOAD]      fait_ventes       : 36203 lignes chargees
PIPELINE TERMINE EN 14 secondes
\end{verbatim}
\end{successbox}

\newpage

% ============================================================
\section{Étape 3 — Implémentation PostgreSQL (20 pts)}

\subsection{Architecture en trois schémas}

Le Data Warehouse utilise trois schémas PostgreSQL à responsabilités séparées :

\begin{lstlisting}[style=sql, caption=Création des schémas PostgreSQL]
CREATE SCHEMA IF NOT EXISTS staging_mexora;   -- zone brute / temporaire
CREATE SCHEMA IF NOT EXISTS dwh_mexora;       -- schema en etoile (dimensions + faits)
CREATE SCHEMA IF NOT EXISTS reporting_mexora; -- vues materialisees analytiques
\end{lstlisting}

\subsection{Dimensions et table de faits}

\begin{lstlisting}[style=sql, caption=DIM\_PRODUIT — avec colonnes SCD Type 2]
CREATE TABLE dwh_mexora.dim_produit (
    id_produit_sk  SERIAL PRIMARY KEY,
    id_produit_nk  VARCHAR(20)  NOT NULL,      -- cle naturelle source
    nom_produit    VARCHAR(200) NOT NULL,
    categorie      VARCHAR(100),
    sous_categorie VARCHAR(100),
    marque         VARCHAR(100),
    fournisseur    VARCHAR(100),
    prix_standard  DECIMAL(10,2),
    origine_pays   VARCHAR(50),
    -- Colonnes SCD Type 2
    date_debut  DATE    NOT NULL DEFAULT CURRENT_DATE,
    date_fin    DATE    NOT NULL DEFAULT '9999-12-31',
    est_actif   BOOLEAN NOT NULL DEFAULT TRUE,
    CONSTRAINT uq_dim_produit_version UNIQUE (id_produit_nk, date_debut)
);
\end{lstlisting}

\begin{lstlisting}[style=sql, caption=FAIT\_VENTES — clés étrangères + contraintes CHECK]
CREATE TABLE dwh_mexora.fait_ventes (
    id_vente    BIGSERIAL PRIMARY KEY,
    id_commande VARCHAR(30)  NOT NULL,
    id_date     INTEGER      NOT NULL
                REFERENCES dwh_mexora.dim_temps(id_date),
    id_produit  INTEGER      NOT NULL
                REFERENCES dwh_mexora.dim_produit(id_produit_sk),
    id_client   INTEGER      NOT NULL
                REFERENCES dwh_mexora.dim_client(id_client_sk),
    id_region   INTEGER      NOT NULL
                REFERENCES dwh_mexora.dim_region(id_region),
    id_livreur  INTEGER
                REFERENCES dwh_mexora.dim_livreur(id_livreur),
    -- Mesures additives
    quantite_vendue       INTEGER      NOT NULL CHECK (quantite_vendue > 0),
    montant_ht            DECIMAL(12,2) NOT NULL,
    montant_ttc           DECIMAL(12,2) NOT NULL,  -- TVA 20 %
    -- Mesure semi-additive
    delai_livraison_jours SMALLINT,
    -- Mesure non-additive
    remise_pct            DECIMAL(5,2)  DEFAULT 0,
    statut_commande       VARCHAR(20)
        CHECK (statut_commande IN ('livre','annule','en_cours','retourne'))
);
\end{lstlisting}

\subsection{Indexation}

\begin{lstlisting}[style=sql, caption=Index simples et composites sur FAIT\_VENTES]
-- Index simples sur les cles etrangeres (jointures analytiques)
CREATE INDEX idx_fv_date    ON dwh_mexora.fait_ventes(id_date);
CREATE INDEX idx_fv_produit ON dwh_mexora.fait_ventes(id_produit);
CREATE INDEX idx_fv_client  ON dwh_mexora.fait_ventes(id_client);
CREATE INDEX idx_fv_region  ON dwh_mexora.fait_ventes(id_region);
CREATE INDEX idx_fv_livreur ON dwh_mexora.fait_ventes(id_livreur);

-- Index composite couvrant (requetes analytiques frequentes)
CREATE INDEX idx_fv_date_region
    ON dwh_mexora.fait_ventes(id_date, id_region)
    INCLUDE (montant_ttc, quantite_vendue);

-- Index partiel sur les commandes livrees (filtre dominant)
CREATE INDEX idx_fv_statut_livre
    ON dwh_mexora.fait_ventes(statut_commande)
    WHERE statut_commande = 'livre';
\end{lstlisting}

\subsection{Trois vues matérialisées}

\subsubsection{Vue 1 — \texttt{mv\_ca\_mensuel} (914 lignes)}

\begin{lstlisting}[style=sql, caption=Vue matérialisée — CA mensuel par région et catégorie]
CREATE MATERIALIZED VIEW reporting_mexora.mv_ca_mensuel AS
SELECT t.annee, t.mois, t.libelle_mois, t.periode_ramadan,
       r.region_admin, r.zone_geo, p.categorie,
       SUM(f.montant_ttc)            AS ca_ttc,
       COUNT(DISTINCT f.id_commande) AS nb_commandes,
       SUM(f.quantite_vendue)        AS volume_vendu,
       ROUND(SUM(f.montant_ttc)
         / NULLIF(COUNT(DISTINCT f.id_commande), 0), 2)  AS panier_moyen,
       COUNT(DISTINCT f.id_client)   AS nb_clients_actifs
FROM dwh_mexora.fait_ventes f
JOIN dwh_mexora.dim_temps   t ON f.id_date    = t.id_date
JOIN dwh_mexora.dim_region  r ON f.id_region  = r.id_region
JOIN dwh_mexora.dim_produit p ON f.id_produit = p.id_produit_sk
WHERE f.statut_commande = 'livre'
GROUP BY t.annee, t.mois, t.libelle_mois, t.periode_ramadan,
         r.region_admin, r.zone_geo, p.categorie
WITH DATA;
CREATE INDEX ON reporting_mexora.mv_ca_mensuel(annee, mois);
CREATE INDEX ON reporting_mexora.mv_ca_mensuel(region_admin);
\end{lstlisting}

\subsubsection{Vue 2 — \texttt{mv\_top\_produits} (3\,013 lignes)}

Classement par CA avec fonction fenêtre \kw{RANK()} intra-catégorie et trimestre :

\begin{lstlisting}[style=sql, caption=Vue matérialisée — Top produits par trimestre]
RANK() OVER (
    PARTITION BY t.annee, t.trimestre, p.categorie
    ORDER BY SUM(f.montant_ttc) DESC
) AS rang_dans_categorie
\end{lstlisting}

\subsubsection{Vue 3 — \texttt{mv\_performance\_livreurs} (215 lignes)}

Taux de retard (délai $> 3$ jours) par livreur, zone et mois :

\begin{lstlisting}[style=sql, caption=Vue matérialisée — Taux de retard livreurs]
ROUND(
    COUNT(*) FILTER (WHERE f.delai_livraison_jours > 3) * 100.0
    / NULLIF(COUNT(*), 0), 2
) AS taux_retard_pct
\end{lstlisting}

\subsection{Vérification d'intégrité référentielle}

\begin{lstlisting}[style=sql, caption=Résultats complets de check\_integrity.sql]
-- Volume par table
    objet    | lignes
-------------+--------
 dim_temps   |   2557
 dim_produit |    300
 dim_client  |    898
 dim_region  |     12
 dim_livreur |      6
 fait_ventes |  36203

-- Controles d'integrite referentielle
      controle       | anomalies
---------------------+-----------
 faits_sans_date     |         0
 faits_sans_produit  |         0
 faits_sans_client   |         0
 faits_sans_region   |         0
 quantites_invalides |         0
 montants_negatifs   |         0
\end{lstlisting}

\begin{successbox}[title=Intégrité référentielle : 0 anomalie sur 36\,203 faits]
Toutes les contraintes \kw{CHECK} et \kw{FOREIGN KEY} sont respectées.
Aucune ligne orpheline, aucune quantité négative, aucun montant négatif.
\end{successbox}

\newpage

% ============================================================
\section{Étape 4 — Dashboard de visualisation (15 pts)}

\subsection{Questions analytiques obligatoires}

\subsubsection{Q1 — Évolution du CA par région (sur 12 mois, comparaison N$-$1)}

\begin{lstlisting}[style=sql, caption=KPI — CA mensuel avec évolution N-1]
WITH ca_mensuel AS (
    SELECT annee, mois, region_admin, SUM(ca_ttc) AS ca_total
    FROM reporting_mexora.mv_ca_mensuel
    GROUP BY annee, mois, region_admin
)
SELECT region_admin,
       ca_total                                               AS ca_actuel,
       LAG(ca_total) OVER (PARTITION BY region_admin
                           ORDER BY annee, mois)             AS ca_precedent,
       ROUND((ca_total - LAG(ca_total) OVER (PARTITION BY
              region_admin ORDER BY annee, mois))
             / NULLIF(LAG(ca_total) OVER (PARTITION BY
               region_admin ORDER BY annee, mois), 0)*100,2) AS evolution_pct
FROM ca_mensuel
ORDER BY annee DESC, mois DESC;
\end{lstlisting}

\begin{resultatbox}[title=Résultat]
Tanger \val{7\,929\,134 MAD} ·
Casablanca \val{7\,793\,512 MAD} ·
Rabat \val{6\,035\,814 MAD} ·
Restant (9 villes) \val{23\,826\,835 MAD}
\end{resultatbox}

\subsubsection{Q2 — Top 10 produits trimestriels à Tanger}

Filtrage sur \kw{dim\_region.ville = 'Tanger'} appliqué à \kw{mv\_top\_produits}
avec filtres interactifs trimestre et année (Metabase / Streamlit).

\subsubsection{Q3 — Panier moyen par segment client}

\begin{lstlisting}[style=sql, caption=Panier moyen et part de CA par segment]
SELECT c.segment_client,
       COUNT(DISTINCT f.id_vente)                          AS nb_commandes,
       ROUND(SUM(f.montant_ttc)
             / COUNT(DISTINCT f.id_vente)::decimal, 2)    AS panier_moyen,
       ROUND(SUM(f.montant_ttc) * 100.0
             / SUM(SUM(f.montant_ttc)) OVER (), 2)        AS pct_ca
FROM dwh_mexora.fait_ventes f
JOIN dwh_mexora.dim_client c ON f.id_client = c.id_client_sk
WHERE f.statut_commande = 'livre' AND c.est_actif = TRUE
GROUP BY c.segment_client
ORDER BY panier_moyen DESC;
\end{lstlisting}

\subsubsection{Q4 — Taux de retour par catégorie (avec alertes)}

\begin{lstlisting}[style=sql, caption=Taux de retour — code couleur dans le dashboard]
SELECT p.categorie,
       COUNT(*) FILTER (WHERE f.statut_commande = 'retourne') AS nb_retours,
       COUNT(*)                                               AS nb_total,
       ROUND(COUNT(*) FILTER (WHERE f.statut_commande = 'retourne')
             * 100.0 / NULLIF(COUNT(*), 0), 2)               AS taux_retour_pct
FROM dwh_mexora.fait_ventes f
JOIN dwh_mexora.dim_produit p ON f.id_produit = p.id_produit_sk
GROUP BY p.categorie
ORDER BY taux_retour_pct DESC;
\end{lstlisting}

\begin{resultatbox}[title=Résultat]
Taux global \val{7,87\,\%} — \badge{mexred}{alerte rouge} activée sur
toutes les catégories (seuil CDC : 5\,\%).
Action prioritaire recommandée : audit de la politique de retour.
\end{resultatbox}

\subsubsection{Q5 — Effet Ramadan sur les ventes alimentaires}

\begin{lstlisting}[style=sql, caption=Comparaison Ramadan vs hors Ramadan]
SELECT t.periode_ramadan,
       ROUND(AVG(ca_ttc), 2) AS ca_moyen_mensuel,
       SUM(volume_vendu)     AS volume_total
FROM reporting_mexora.mv_ca_mensuel
WHERE categorie = 'Alimentation'
  AND annee >= 2022
GROUP BY t.periode_ramadan;
\end{lstlisting}

\begin{resultatbox}[title=Résultat]
Indice de performance Ramadan = \val{91,2} (base 100 hors Ramadan).
Les données synthétiques ne reproduisent pas l'effet saisonnier réel ;
l'indicateur est conçu pour fonctionner avec des données transactionnelles réelles.
\end{resultatbox}

\subsection{Version web déployable — Streamlit}

\textbf{Metabase} est conservé comme outil BI local conforme au cahier de charge (livrable L7).
En complément, \textbf{Streamlit} a été ajouté comme interface web déployable pour faciliter
la démonstration du projet en ligne.

\begin{infobox}[title=Architecture Streamlit — découplage du Data Warehouse]
\begin{itemize}
  \item Les données sont exportées depuis PostgreSQL via
    \texttt{scripts/export\_streamlit\_data.py} vers \texttt{dashboard/data/*.csv}.
  \item L'application \texttt{dashboard/streamlit\_app.py} lit uniquement ces CSV :
    \textbf{aucune connexion PostgreSQL n'est requise} sur Streamlit Cloud.
  \item Le dashboard couvre 5 pages : Vue générale, Analyse régionale, Clients,
    Produits, Retours \& Ramadan.
  \item Lien déployé :
    \url{https://pipeline-etl-data-warehouse.streamlit.app/}
\end{itemize}
\end{infobox}

\begin{lstlisting}[style=bash, caption=Workflow complet — export et déploiement Streamlit]
# 1. Exporter les donnees depuis PostgreSQL
python scripts/export_streamlit_data.py

# 2. Lancer localement (http://localhost:8501)
streamlit run dashboard/streamlit_app.py

# 3. Deployer sur Streamlit Cloud
#    -> share.streamlit.io  |  repo : SoufianeZaari/Pipeline-ETL-Data-Warehouse
#    -> main file           :  dashboard/streamlit_app.py
\end{lstlisting}

\subsection{Synthèse des insights métier}

\rowcolors{2}{rowalt}{white}
\begin{table}[H]
\centering\small
\begin{tabular}{p{3.5cm} p{9.5cm}}
\toprule
\rowcolor{mexlightblue}
\textbf{Indicateur} & \textbf{Insight et recommandation} \\
\midrule
CA total &
  \val{45\,585\,295 MAD} sur le périmètre chargé \\[3pt]
Top région &
  Tanger $+$1,7\,\% vs Casablanca — consolider le marché local tangérois \\[3pt]
Top catégorie &
  Electronics 82\,\% du CA — sécuriser les stocks et les fournisseurs \\[3pt]
Taux de retour &
  7,87\,\% (seuil : 5\,\%) — auditer la politique retour et les fiches produit \\[3pt]
Délai livraison &
  14,99 jours en moyenne — objectif $<\,7$ jours via SLA transporteurs \\[3pt]
Segment Gold &
  Panier moyen le plus élevé — programme fidélité VIP à prioriser \\
\bottomrule
\end{tabular}
\caption{Insights métier identifiés dans le Data Warehouse Mexora}
\end{table}
\rowcolors{0}{}{}

\newpage

% ============================================================
\section{Conclusion}

Ce mini-projet démontre la mise en œuvre complète d'un système décisionnel d'entreprise,
de l'analyse des besoins jusqu'au dashboard déployé, en respectant rigoureusement les
quatre étapes du cahier de charge.

\begin{infobox}[title=Bilan technique par étape]
\begin{tabular}{@{}p{1.8cm} p{11.5cm}@{}}
\textbf{Étape 1} &
  Modélisation en étoile rigoureuse (5 dimensions), granularité justifiée, structure
  SCD Type\,2 préparée via \kw{date\_debut}, \kw{date\_fin}, \kw{est\_actif} sur
  DIM\_PRODUIT et DIM\_CLIENT (historisation multi-version activable dans les évolutions
  futures du DWH). \\[5pt]
\textbf{Étape 2} &
  Pipeline Python structuré \texttt{mexora\_etl/}, 14 règles de transformation
  documentées et loguées, segmentation client automatisée, 14 secondes d'exécution
  pour 50\,000 lignes source. \\[5pt]
\textbf{Étape 3} &
  PostgreSQL 16 avec 3 schémas dédiés, 36\,203 faits chargés, 3 vues matérialisées
  fonctionnelles, 0 anomalie sur 6 contrôles d'intégrité référentielle. \\[5pt]
\textbf{Étape 4} &
  Metabase comme dashboard BI local conforme au cahier de charge, 5 questions
  analytiques avec KPIs et alertes visuelles. Complément : application Streamlit
  déployée en ligne, autonome via des exports CSV du Data Warehouse.
\end{tabular}
\end{infobox}

\subsection*{Hypothèses et choix assumés}

\begin{enumerate}
  \item \textbf{TVA 20\,\%} appliquée uniformément
    (\kw{montant\_ttc = montant\_ht $\times$ 1.20})
    en l'absence de taux différenciés par catégorie dans les données source.
  \item \textbf{Livreurs fictifs} (L001--L005 + inconnu $-1$) enrichis avec
    des métadonnées statiques, faute de référentiel transporteur dans les sources.
  \item \textbf{Statuts inconnus} exclus de la table de faits (contrainte
    \kw{CHECK} PostgreSQL) et tracés dans les logs ETL.
  \item \textbf{Instance PostgreSQL locale} sur le port 5433
    (socket \texttt{/tmp/mexora\_pg\_socket}) pour contournement des droits
    administrateur sur la machine académique.
  \item \textbf{Segmentation Streamlit} recalculée par percentiles (P30/P70) car le CA
    synthétique par client dépasse largement les seuils fixes (67\,K--668\,K MAD),
    rendant la classification fixe mono-segment. Le schéma DWH n'est pas modifié.
\end{enumerate}

\vspace{10pt}
\begin{center}
{\color{mexblue}\rule{0.6\linewidth}{1.5pt}}\\[6pt]
{\small\itshape
  Rapport — Miniprojet 1 · Mexora Analytics · Soufiane Zaari · 2025--2026}\\[3pt]
{\small\url{https://github.com/SoufianeZaari/Pipeline-ETL-Data-Warehouse}}
\end{center}

\end{document}
